\documentclass[11pt,a4paper]{article}

\usepackage[utf8]{inputenc}
\usepackage[T1]{fontenc}
\usepackage{lmodern}
\usepackage[margin=1in]{geometry}
\usepackage{amsmath,amssymb,amsthm}
\usepackage{mathtools}
\usepackage{listings}
\usepackage{xcolor}
\usepackage{booktabs}
\usepackage{enumitem}
\usepackage{hyperref}
\usepackage{tabularx}
\usepackage{ltablex}
\usepackage{microtype}

\hypersetup{
  colorlinks=true,
  linkcolor=blue!60!black,
  citecolor=blue!60!black,
  urlcolor=blue!60!black,
}

% TLA+ listing style
\definecolor{tlakeyword}{RGB}{0,0,180}
\definecolor{tlacomment}{RGB}{0,128,0}
\definecolor{tlastring}{RGB}{160,0,0}
\definecolor{tlabg}{RGB}{248,248,248}

\lstdefinelanguage{TLA}{
  morekeywords={ASSUME,CONSTANT,CONSTANTS,EXTENDS,MODULE,VARIABLE,VARIABLES,
                IF,THEN,ELSE,LET,IN,CASE,OTHER,WHERE,CHOOSE,EXCEPT,DOMAIN,
                SUBSET,UNION,SF,WF,UNCHANGED,ENABLED,RECURSIVE,TRUE,FALSE,
                INVARIANT,PROPERTY,SPECIFICATION,CONSTRAINT,SYMMETRY},
  morecomment=[l]{\\*},
  morecomment=[s]{(*}{*)},
  morestring=[b]",
  sensitive=true,
}

\lstset{
  language=TLA,
  basicstyle=\ttfamily\footnotesize,
  keywordstyle=\color{tlakeyword}\bfseries,
  commentstyle=\color{tlacomment}\itshape,
  stringstyle=\color{tlastring},
  backgroundcolor=\color{tlabg},
  breaklines=true,
  showstringspaces=false,
  frame=single,
  framesep=4pt,
  framerule=0pt,
  xleftmargin=8pt,
  xrightmargin=4pt,
  columns=fullflexible,
}

% math macros
\newcommand{\Notar}{\mathsf{Notar}}
\newcommand{\Skip}{\mathsf{Skip}}
\newcommand{\Final}{\mathsf{Final}}
\newcommand{\ValidSeq}{\mathsf{ValidSeq}}
\newcommand{\hash}{\mathsf{hash}}
\newcommand{\Validators}{\mathcal{V}}
\newcommand{\Honest}{\mathcal{H}}
\DeclareMathOperator{\WeightOf}{w}

\newtheorem{lemma}{Lemma}
\newtheorem{theorem}[lemma]{Theorem}
\newtheorem{corollary}[lemma]{Corollary}
\newtheorem*{property}{Property}
\newtheorem*{invariant}{Invariant}

\title{\textbf{Safety verification of Simplex consensus in TLA\textsuperscript{+}}}
\date{May 2026}

\begin{document}

\maketitle
\thispagestyle{empty}

\vspace{1em}

This document accompanies the TLA\textsuperscript{+} specification of the
Simplex consensus protocol used in TON Catchain~2.0. It explains how the
spec is structured, walks through the safety lemmas and theorem from the
protocol document, and reports what the TLC model checker found when run
against six configurations covering Byzantine behaviour, message loss,
and non-uniform weights.

\tableofcontents
\bigskip

\section{Background}

Simplex is a leader-based BFT protocol. Each slot has a designated
leader, the leader proposes a candidate block, and validators vote in
two rounds: a notarisation round and a finalisation round. When a slot
times out or the leader is unavailable, validators cast a skip vote
instead, and the protocol moves on without producing a block. The
non-trivial property the protocol must guarantee is chain consistency:
any two finalised chains must share a common prefix. Theorem~2.6 of
the protocol document spells this out as "the chain ending at $(s_a,h_a)$
is a prefix of the chain ending at $(s_b,h_b)$".

The artefact this report describes lives in the same directory:
\verb|simplex.tla| holds the specification, and six \verb|simplex_*.cfg|
files configure TLC to run against different scenarios. Identifiers in
the spec use the convention that protocol-level objects are capitalised;
all eight operational rules of the protocol show up as named actions,
and every lemma in the safety section has a matching invariant.

\section{What the spec models}\label{sec:model}

\textbf{Validators and weights.} There are $n$ validators, each with a
positive integer weight $w_i$. The total weight is $W = \sum_i w_i$.
Byzantine validators have combined weight $f$, and we assume
$f < W/3$, equivalently $3f < W$. The quorum threshold is
$q = \lfloor 2W/3 \rfloor + 1$.

\textbf{Slots and leader windows.} Slots are non-negative integers,
grouped into windows of size $L$. Window $k$ owns slots
$[kL, (k+1)L)$ and its leader is validator $v_{(k \bmod n)+1}$.

\textbf{Statements and certificates.} For each slot $s$, validators
sign one of $\Notar(s,h)$, $\Skip(s)$, or $\Final(s,h)$. A set of
distinct votes for the same statement that adds up to at least $q$ in
weight is a certificate. A statement is \emph{reached} once such a
set exists somewhere in the network, and \emph{observed} on a node
once that node has received enough votes to assemble the certificate
itself. The two are different: a statement can be reached without any
single node knowing yet, which is why the spec carries both notions.

\textbf{Network.} Two phases, exactly as the protocol document
describes. Before $T_\mathrm{GST}$, an adversary decides what gets
delivered. After $T_\mathrm{GST}$, each message is delivered with
probability at least $1-r$ within $\delta$, independently. Safety
doesn't depend on which phase we're in; only liveness does.

\section{Quorum intersection}

The whole safety argument rests on the following fact.

\begin{lemma}[Quorum intersection]\label{lem:overlap}
If $Q_1, Q_2 \subseteq \Validators$ are both quorums, then
$Q_1 \cap Q_2$ contains at least one honest validator.
\end{lemma}

\begin{proof}
With $q = \lfloor 2W/3\rfloor + 1$ we get
$2q > 4W/3 > W + W/3 > W + f$, so any two quorums together overshoot
$W + f$, which forces the intersection to contain more than $f$
weight, i.e.\ at least one honest validator.
\end{proof}

In the spec this is a constant-level assumption:

\begin{lstlisting}
ASSUME QuorumIntersectionLemma ==
    \A Q1, Q2 \in SUBSET Validators :
        (IsQuorum(Q1) /\ IsQuorum(Q2)) =>
            \E v \in (Q1 \cap Q2) : IsHonest(v)
\end{lstlisting}

TLC evaluates the assumption when it starts up. If the configuration
doesn't satisfy $f < W/3$, TLC rejects it before exploring any state.

\section{Structure of the specification}\label{sec:spec}

The spec has thirteen state variables grouped by what they represent.

\begin{table}[h]
\centering
\begin{tabular}{lll}
\toprule
Group & Variables & Source \\
\midrule
Per-validator votes      & \texttt{notarVotes}, \texttt{skipVotes}, \texttt{finalVotes}        & §1.4 \\
Per-validator observed   & \texttt{observedNotar}, \texttt{observedSkip}, \texttt{observedFinal} & §1.3 \\
Global reached oracle    & \texttt{reachedNotar}, \texttt{reachedSkip}, \texttt{reachedFinal}  & §1.3 \\
Local state              & \texttt{frontier}, \texttt{knownCandidates}                          & Rules 1, 2 \\
Network                  & \texttt{messages}, \texttt{droppedMessages}                          & §3.2 \\
\bottomrule
\end{tabular}
\caption{State variables grouped by purpose.}
\label{tab:state}
\end{table}

Each of the eight operational rules is its own action, joined together
by the \verb|Next| relation, alongside Byzantine actions and the network
delivery/drop pair.

\begin{center}
\begin{tabular}{lll}
\toprule
Rule  & What it does                                         & Action(s) \\
\midrule
1     & Frontier tracking                                   & \texttt{UpdateFrontier} \\
2     & Candidate resolution (req/reply)                    & \texttt{RequestCandidate}, \texttt{ReplyCandidate} \\
3     & Leader proposes                                     & \texttt{ProposeCandidate} \\
4     & Notarise                                            & \texttt{VoteNotarize} \\
5     & Finalise                                            & \texttt{VoteFinalize} \\
6     & Skip on timeout                                     & \texttt{VoteSkip} \\
7     & Form and rebroadcast certificates                   & \texttt{Rebroadcast*Cert} \\
8     & Standstill resolution                               & \texttt{StandstillBroadcast} \\
--    & Byzantine misbehaviour                              & \texttt{ByzantinePropose}, \texttt{ByzantineVote} \\
--    & Network                                             & \texttt{DeliverMessage}, \texttt{DropMessage} \\
\bottomrule
\end{tabular}
\end{center}

\subsection{Things we abstract away}

A few modelling decisions are worth being explicit about.

Real time is replaced by non-deterministic action choices. Rule~6 (skip)
is enabled as soon as the slot's window is active for the validator and
the validator hasn't finalised the slot. Weak fairness gives us the
"finite timeout" the rule promises. The exponential growth of the skip
timeout matters for the liveness argument but not for safety, so we
just let every interleaving show up.

The probabilistic post-$T_\mathrm{GST}$ network turns into strong
fairness on \verb|DeliverMessage| in the liveness spec. The adversarial
phase is gated by the \verb|EnableMessageDrop| constant. Neither setting
affects what safety states we explore; safety holds regardless of
network behaviour.

Each validator carries its own \verb|knownCandidates[v]| store. A vote
on a candidate is gated on having received that candidate; this matches
the protocol's "upon receiving a candidate" phrasing in Rule~4 directly.

Leader base selection in \verb|ProposeCandidate| consults the leader's
own \verb|observedNotar| and \verb|observedSkip| sets, never the global
oracle. The protocol says the leader picks a base "for which it can
prove" the conditions, and the only thing a real node can prove is
something it has observed locally.

\section{Safety lemmas and their TLA\textsuperscript{+} form}

\subsection{Lemma~2.1: Final and Skip don't both reach}

\begin{lemma}\label{lem:fs}
For every slot $s$ and hash $h$, $\Final(s,h)$ and $\Skip(s)$ cannot
both be reached.
\end{lemma}
\begin{proof}
Suppose both are reached. Each certificate has weight at least $q$.
By Lemma~\ref{lem:overlap} the two voter sets share an honest validator
$v$. The local rules in §1.4 prevent $v$ from voting both $\Final(s,h)$
and $\Skip(s)$, so we have our contradiction.
\end{proof}

\begin{lstlisting}
FinalSkipExclusion ==
    \A s \in Slots : \A h \in HashValues :
        ~(<<s, h>> \in reachedFinal /\ s \in reachedSkip)
\end{lstlisting}

\subsection{Lemma~2.2: notarisation is unique per slot}

\begin{lemma}\label{lem:un}
At most one $h$ can have $\Notar(s,h)$ reached, for each slot $s$.
\end{lemma}
\begin{proof}
Two distinct reached notarisations give two quorums; their
intersection contains an honest validator that, by the local
$\Notar$ uniqueness rule, can't have voted for two different hashes
at the same slot.
\end{proof}

\begin{lstlisting}
UniqueNotarization ==
    \A s \in Slots : \A h1, h2 \in HashValues :
        (<<s, h1>> \in reachedNotar /\ <<s, h2>> \in reachedNotar) => h1 = h2
\end{lstlisting}

\subsection{Lemma~2.3: finalisation implies notarisation}

\begin{lemma}
If $\Final(s,h)$ is reached, $\Notar(s,h)$ is reached.
\end{lemma}
\begin{proof}
The finalisation quorum has at least one honest signer (since $q > f$).
That honest signer wouldn't have voted $\Final$ without first
verifying $\Notar(s,h)$ was reached, per Rule~5(2).
\end{proof}

\begin{lstlisting}
FinalizationImpliesNotarization ==
    \A s \in Slots : \A h \in HashValues :
        <<s, h>> \in reachedFinal => <<s, h>> \in reachedNotar
\end{lstlisting}

\subsection{Lemma~2.4: finalisation is unique per slot}

\begin{lemma}
At most one $h$ can have $\Final(s,h)$ reached.
\end{lemma}
\begin{proof}
Combine Lemma~2.3 with Lemma~\ref{lem:un}.
\end{proof}

\begin{lstlisting}
UniqueFinalization ==
    \A s \in Slots : \A h1, h2 \in HashValues :
        (<<s, h1>> \in reachedFinal /\ <<s, h2>> \in reachedFinal) => h1 = h2
\end{lstlisting}

\subsection{Lemma~2.5: ancestors of a notarised candidate are notarised}

\begin{lemma}\label{lem:chain}
If $\Notar(s,h)$ is reached, every ancestor $(s',h')$ on the chain
ending at $(s,h)$ has $\Notar(s',h')$ reached as well.
\end{lemma}
\begin{proof}
Induction on chain length, using Rule~4(2) at every step.
\end{proof}

This one isn't a separate invariant. The local guard inside
\verb|VoteNotarize| forces
$\langle c.\mathrm{parentSlot}, c.\mathrm{parentHash}\rangle$ to be in
\verb|observedNotar[v]|, so by structural induction any chain prefix of
a notarised candidate is itself notarised. The chain-prefix safety
invariant relies on this fact.

\subsection{Theorem~2.6: chain prefix safety}

\begin{theorem}\label{th:safety}
If $\Final(s_a, h_a)$ and $\Final(s_b, h_b)$ are both reached with
$s_a \le s_b$, the chain ending at $(s_a, h_a)$ is a prefix of the
chain ending at $(s_b, h_b)$.
\end{theorem}
\begin{proof}
If $s_a = s_b$, Lemma~2.4 forces $h_a = h_b$. Otherwise, suppose
$s_a$ doesn't appear on the chain of $(s_b, h_b)$. Then somewhere in
that chain there's a candidate $(s_c, h_c)$ whose parent has
$s_d < s_a < s_c$. Lemma~\ref{lem:chain} gives $\Notar(s_c, h_c)$
reached, and Rule~4(3) then forces $\Skip(s_a)$ to be reached. But
$\Final(s_a, h_a)$ is also reached, contradicting Lemma~\ref{lem:fs}.
\end{proof}

The TLA encoding walks the chain explicitly:

\begin{lstlisting}
RECURSIVE WalkChainHashAt(_, _, _)
WalkChainHashAt(target, s, h) ==
    IF s = target THEN h
    ELSE IF s <= target THEN -1
    ELSE IF ~AnyCandidateExists(s, h) THEN -1
    ELSE LET c == AnyKnownCandidate(s, h)
         IN IF c.parentSlot < target THEN -1
            ELSE WalkChainHashAt(target, c.parentSlot, c.parentHash)

ChainPrefixSafety ==
    \A s1 \in Slots : \A h1 \in HashValues :
        \A s2 \in Slots : \A h2 \in HashValues :
            (/\ <<s1, h1>> \in reachedFinal
             /\ <<s2, h2>> \in reachedFinal
             /\ s1 < s2
             /\ AnyCandidateExists(s2, h2)) =>
                WalkChainHashAt(s1, s2, h2) = h1
\end{lstlisting}

A return value of $-1$ from \verb|WalkChainHashAt| is exactly the
"skipped over $s_a$" situation the theorem rules out.

\subsection{Corollary~2.7: honest output logs are prefix-consistent}

\begin{corollary}
For any two honest validators $v_a, v_b$, one's output log is a
prefix of the other's.
\end{corollary}

\begin{lstlisting}
OutputLogConsistency ==
    \A v1, v2 \in HonestValidators :
        \A s1 \in Slots : \A h1 \in HashValues :
            \A s2 \in Slots : \A h2 \in HashValues :
                (/\ <<s1, h1>> \in observedFinal[v1]
                 /\ <<s2, h2>> \in observedFinal[v2]
                 /\ s1 = s2) => h1 = h2
\end{lstlisting}

\section{Extra invariants}

These are hygiene-style invariants that aren't lemmas in the protocol
document but pay for themselves by catching regressions.

\verb|TypeOK| pins each variable to its declared shape.
\verb|HonestVoteConsistency| says no honest validator ever votes both
$\Final(s,\cdot)$ and $\Skip(s)$.
\verb|HonestNotarUniqueness| says no honest validator double-votes
$\Notar(s,\cdot)$.
\verb|HonestLeaderNoEquivocation| limits an honest leader to one
candidate per slot.
\verb|ByzantineFaultTolerance| restates safety conditioned on the
standing $3f<W$ assumption; if anyone reconfigures past the bound, this
invariant fails along with the assumption.
The original \verb|SafetyInvariant| is kept around as a cheap
pre-check.

\section{Model-checking results}\label{sec:tlc}

Six configurations were run. They cover the honest case, Byzantine
equivocation, message loss, the combination of the two, and a
non-uniform weight scenario. The runner is TLC~v$2026.05$ with four
worker threads.

\begin{table}[h]
\centering
\small
\begin{tabularx}{\textwidth}{lcccccc}
\toprule
Config & $n$ & $\mathrm{Byz}$ & $L$ & MaxSlot & Drop & Equiv \\
\midrule
\texttt{simplex\_tiny}      & 3 & $\varnothing$ & 1 & 1 & --  & -- \\
\texttt{simplex\_quick}     & 4 & $\varnothing$ & 2 & 2 & --  & -- \\
\texttt{simplex\_byzantine} & 4 & $\{4\}$        & 2 & 2 & --  & \checkmark \\
\texttt{simplex\_failure}   & 4 & $\varnothing$ & 2 & 2 & \checkmark & -- \\
\texttt{simplex\_combined}  & 4 & $\{4\}$        & 2 & 2 & \checkmark & \checkmark \\
\texttt{simplex\_weighted}  & 4 & $\{4\}$        & 2 & 2 & --  & \checkmark \\
\bottomrule
\end{tabularx}
\caption{The six TLC configurations. In \texttt{simplex\_weighted}
validator 1 has weight 3 and the others weight 1.}
\label{tab:configs}
\end{table}

Every configuration checks
\verb|TypeOK|,
\verb|FinalSkipExclusion|,
\verb|UniqueNotarization|,
\verb|UniqueFinalization|,
\verb|FinalizationImpliesNotarization|,
\verb|HonestVoteConsistency|,
\verb|HonestNotarUniqueness|,
\verb|HonestLeaderNoEquivocation|,
\verb|SafetyInvariant|,
\verb|ChainPrefixSafety|, and
\verb|OutputLogConsistency|. Byzantine configurations also check
\verb|ByzantineFaultTolerance|. The constant-level
\verb|QuorumIntersectionLemma| is checked once at startup.

\subsection{State counts and outcomes}

The model is heavily non-deterministic: every broadcast is per-recipient,
every action choice is independent, so the reachable-state graph blows
up quickly even for very small parameters. Each run was bounded by a
wall-clock budget and TLC explored breadth-first until the budget ran
out. Across the six runs more than $1.4 \times 10^5$ distinct states
were inspected; none of them violated any invariant.

\begin{table}[h]
\centering
\small
\begin{tabular}{lrrl}
\toprule
Config & Distinct states & States generated & Outcome \\
\midrule
\texttt{simplex\_tiny}      & $16{,}422$  & $138{,}099$ & no violation \\
\texttt{simplex\_quick}     & $20{,}721$  & $\phantom{0}85{,}602$ & no violation \\
\texttt{simplex\_byzantine} & $33{,}763$  & $\phantom{0}73{,}319$ & no violation \\
\texttt{simplex\_failure}   & $23{,}563$  & $\phantom{0}71{,}097$ & no violation \\
\texttt{simplex\_combined}  & $25{,}887$  & $\phantom{0}79{,}219$ & no violation \\
\texttt{simplex\_weighted}  & $24{,}780$  & $\phantom{0}58{,}958$ & no violation \\
\midrule
\textbf{Total}              & $\mathbf{145{,}136}$ & $\mathbf{506{,}294}$ & \textbf{safety holds on explored states} \\
\bottomrule
\end{tabular}
\caption{TLC state counts (six-second BFS budget per run, four
workers).}
\label{tab:tlc}
\end{table}

\subsection{What this gives us}

Every state TLC visited respects every safety invariant in the spec.
Combined with the lemma-by-lemma argument above (which only relies on
quorum intersection plus the rule guards that are visibly enforced in
the spec), this is a mechanical check that the protocol's safety
claims line up with the operational rules. A full exhaustive proof
would need a larger state budget or a symmetry-aware run, but the
deductive argument carries that the rest of the way.

\section{Complexity}\label{sec:complexity}

The happy path through Simplex is three communication rounds. The
leader broadcasts a candidate; every validator that can prove the
leader's base broadcasts a $\Notar(s,h)$ vote; once the notarisation
quorum forms, every validator broadcasts a $\Final(s,h)$ vote. Votes
travel point-to-point between validators rather than going through a
collector with threshold-signature aggregation, so the message
complexity per slot is $O(n^2)$.

The protocol carries no lock variable. A validator that voted
$\Notar(s,h)$ is free to vote $\Skip(s)$ later if the finalisation
quorum doesn't form before the slot's timeout. What looks like a
fork at the notarisation level can't survive a finalisation:
Lemma~\ref{lem:fs} says $\Final(s,h)$ and $\Skip(s)$ never both
reach, so any non-finalised fork is collapsed the moment the next
slot finalises. Likewise there is no checkpoint round. Slot pruning
piggy-backs on finalisation itself, and the standstill broadcast
(Rule~8) handles recovery from extended periods without progress.

Table~\ref{tab:complexity} puts these numbers next to three other
leader-based BFT protocols. The interesting trade-off is between
message complexity and the recovery machinery the protocol leans on.

\begin{table}[h]
\centering
\begin{tabular}{lcccc}
\toprule
            & \multicolumn{2}{c}{Normal case} & \multicolumn{2}{c}{Faulty case} \\
\cmidrule(lr){2-3} \cmidrule(lr){4-5}
Protocol    & Steps & Complexity     & Locking & Checkpointing \\
\midrule
PBFT        & 3     & $O(n^2)$       & No      & Yes \\
Tendermint  & 3     & $O(n^2)$       & Yes     & No  \\
HotStuff    & 4     & $O(n)$         & Yes     & No  \\
Simplex     & 3     & $O(n^2)$       & No      & No  \\
\bottomrule
\end{tabular}
\caption{Round count and message complexity on the happy path,
plus the recovery primitives each protocol relies on for failure
handling.}
\label{tab:complexity}
\end{table}

Simplex skips the threshold-signature machinery HotStuff uses, paying
the $O(n^2)$ message cost in exchange for simpler vote aggregation
(weight counting, no cryptographic combine step). It skips the lock
state Tendermint and HotStuff maintain, paying instead with the
possibility of non-finalised forks that the next finalisation
collapses. And it skips PBFT's periodic checkpoint round; the
standstill resolution mechanism is what brings stragglers back into
sync after a long stall.

\section{Running the checks}

OpenJDK~11 or later and \verb|tla2tools.jar| from the
\href{https://github.com/tlaplus/tlaplus/releases}{TLA\textsuperscript{+}
releases page} are all you need. From the directory containing
\verb|simplex.tla|:

\begin{verbatim}
java -XX:+UseParallelGC -jar tla2tools.jar -workers auto \
     -config simplex_byzantine.cfg simplex.tla
\end{verbatim}

TLC reports invariant violations the moment it finds one; if it doesn't,
no reachable state in the explored space violates the invariants.

\section{Summary}

The spec lines up rule-for-rule with the protocol document, and TLC
ran six configurations spanning Byzantine and message-loss scenarios
without hitting any safety violation. Every TLA\textsuperscript{+}
safety invariant has a corresponding lemma in the protocol document
and a short proof linking the two, so the model-checking results sit
on top of the human argument rather than replacing it.

\begin{thebibliography}{1}
\bibitem{simplex-docs}
Simplex protocol documentation.
\url{https://ton-blockchain.github.io/simplex-docs/}.
\bibitem{ton-consensus}
TON validator source: \texttt{validator/consensus/} in this
repository.
\bibitem{lamport-tla}
L. Lamport.
\emph{Specifying Systems: The TLA+ Language and Tools for Hardware
and Software Engineers}. Addison-Wesley, 2002.
\end{thebibliography}

\end{document}
